\documentclass[11pt,a4paper]{article}
\usepackage{fontspec}
\setmainfont{Nirmala UI}
\usepackage{geometry}
\geometry{a4paper, margin=1in}
\usepackage{hyperref}
\usepackage{color}
\usepackage{enumitem}
\usepackage{titlesec}
\usepackage{fancyhdr}
\usepackage{longtable}

% Professional colors
\definecolor{primary}{rgb}{0.1, 0.3, 0.6}
\definecolor{secondary}{rgb}{0.3, 0.3, 0.3}

\hypersetup{
    colorlinks=true,
    linkcolor=primary,
    filecolor=primary,      
    urlcolor=primary,
    pdftitle={AIGGPA Fieldwork Readiness Dossier},
}

\titleformat{\section}{\large\bfseries\color{primary}}{\thesection}{1em}{}
\titleformat{\subsection}{\normalsize\bfseries\color{secondary}}{\thesubsection}{1em}{}

\pagestyle{fancy}
\fancyhf{}
\rhead{\textcolor{secondary}{AIGGPA Research Project}}
\lhead{\textcolor{secondary}{Fieldwork Readiness Dossier}}
\rfoot{\thepage}

\title{\LARGE \bfseries \color{primary} AIGGPA Research Fieldwork Readiness Dossier \\ \Large Comprehensive Project Overview \& Instrument Repository}
\author{\textbf{Aryan Kori} \\ Intern, AIGGPA}
\date{\today}

\begin{document}

\maketitle
\tableofcontents
\newpage

\section{Executive Summary}
This dossier provides a comprehensive, line-by-line detailed account of the research preparation for the study titled \textbf{"Assessment of Digital Tool Adoption and Its Impact on Efficiency in MP Government Departments"}. 

Every aspect of the workspace has been documented here, including the exact questions asked in the surveys, the purpose of each automation script, the folder structure on Google Drive and GitHub, and the proposed fieldwork itinerary. This document serves as the single source of truth for the project manager to review before approving the commencement of fieldwork.

\section{Project Overview \& Objectives}
The primary objective of this research is to evaluate how effectively digital tools are being adopted by government employees across four key departments in Madhya Pradesh:
\begin{enumerate}[label=\alph*)]
    \item Revenue Department
    \item Rural Development Department
    \item Forest Department
    \item Health Department
\end{enumerate}

The study aims to answer:
\begin{itemize}
    \item What is the actual adoption rate of official portals vs. manual registers?
    \item Do digital tools increase or decrease workload for field staff?
    \item What are the infrastructure gaps (internet, devices) preventing full adoption?
    \item To what extent are personal tools (like WhatsApp) filling the gaps left by official systems?
\end{itemize}

\section{Methodology \& Instrument Design}
A hybrid approach (Digital + Print) has been adopted. All schedules are **bilingual (English and Hindi)** to ensure no respondent is excluded due to language barriers.

\subsection{Folder Structure \& Infrastructure}
The workspace is organized to ensure data integrity and backup across local storage, Google Drive, and GitHub:
\begin{itemize}
    \item \textbf{`scripts/`}: Contains all Python automation scripts.
    \item \textbf{`schedules/pdf/`}: Stores all generated PDF schedules.
    \item \textbf{`schedules/docx/`}: Stores all converted Word files.
    \item \textbf{`schedules/tex/`}: Stores all source LaTeX files.
\end{itemize}

\section{Detailed Script Reference}
Here is the list of all scripts created in the workspace and their specific functions:

\begin{enumerate}
    \item \textbf{\texttt{gen\_printable\_schedule.py}}
    \begin{itemize}
        \item \textbf{Purpose:} Generates the master bilingual questionnaire in LaTeX and compiles it to PDF.
    \end{itemize}
    
    \item \textbf{\texttt{gen\_dept\_schedules.py}}
    \begin{itemize}
        \item \textbf{Purpose:} Generates 4 separate PDF schedules, one for each department.
    \end{itemize}
    
    \item \textbf{\texttt{create\_form.py}}
    \begin{itemize}
        \item \textbf{Purpose:} Automates the creation of the giant Google Form using the \texttt{gog} CLI tool.
    \end{itemize}
    
    \item \textbf{\texttt{hi.py}}
    \begin{itemize}
        \item \textbf{Purpose:} The core translation dictionary. Stores all questions in both English and Hindi.
    \end{itemize}
    
    \item \textbf{\texttt{convert\_pdf2docx\_v2.py}}
    \begin{itemize}
        \item \textbf{Purpose:} Converts the generated PDFs into editable Microsoft Word documents.
    \end{itemize}
    
    \item \textbf{\texttt{upload\_to\_drive.py}}
    \begin{itemize}
        \item \textbf{Purpose:} Automates the uploading of all categorized files to Google Drive.
    \end{itemize}
\end{enumerate}

\newpage
\appendix
\section{Full Survey Instrument (The Questions)}
Here is the exact list of questions configured in the system.

\subsection{Common Questions (Q1 - Q39)}
These questions are asked to ALL respondents across all departments.

\begin{longtable}{|p{1cm}|p{6cm}|p{6cm}|}
\hline
\textbf{No.} & \textbf{English Question} & \textbf{Hindi Translation} \\ \hline
\endfirsthead
\hline
\textbf{No.} & \textbf{English Question} & \textbf{Hindi Translation} \\ \hline
\endhead
\hline
\endfoot
\hline
\endlastfoot
Q1 & Designation / Post: & पदनाम / पद: \\ \hline
Q2 & Job Role / Level: & कार्य भूमिका / स्तर: \\ \hline
Q3 & Age group: & आयु वर्ग: \\ \hline
Q4 & Gender: & लिंग: \\ \hline
Q5 & Years of service: & सेवा के वर्ष: \\ \hline
Q6 & Highest education: & उच्चतम शिक्षा: \\ \hline
Q7 & What digital devices are available at your workstation? & आपके कार्यस्थल पर कौन-से डिजिटल उपकरण उपलब्ध हैं? \\ \hline
Q8 & Do you share your device with other employees? & क्या आप अपना उपकरण अन्य कर्मचारियों के साथ साझा करते हैं? \\ \hline
Q9 & Rate internet connectivity at your office: & अपने कार्यालय में इंटरनेट कनेक्टिविटी का मूल्यांकन करें: \\ \hline
Q10 & How often do you experience internet outages per week? & सप्ताह में कितनी बार इंटरनेट बंद होता है? \\ \hline
Q11 & Is IT helpdesk / technical support available? & क्या IT हेल्पडेस्क / तकनीकी सहायता उपलब्ध है? \\ \hline
Q12 & If yes, how quickly are issues resolved? & यदि हाँ, तो समस्याएँ कितनी जल्दी हल होती हैं? \\ \hline
Q13 & Digital tools help me complete tasks faster than paper. & डिजिटल उपकरण मुझे कागज़ की तुलना में कार्य तेज़ी से पूरा करने में मदद करते हैं। \\ \hline
Q14 & Digital tools improve the quality/accuracy of my work. & डिजिटल उपकरण मेरे काम की गुणवत्ता/सटीकता में सुधार करते हैं। \\ \hline
Q15 & Using digital tools increases my overall productivity. & डिजिटल उपकरणों का उपयोग मेरी समग्र उत्पादकता बढ़ाता है। \\ \hline
Q16 & The digital tools available are well-suited to my actual job tasks. & उपलब्ध डिजिटल उपकरण मेरे वास्तविक कार्यों के लिए उपयुक्त हैं। \\ \hline
Q17 & How difficult do you find digital tools to use? & डिजिटल उपकरणों का उपयोग आपको कितना कठिन लगता है? \\ \hline
Q18 & I am confident in my ability to use digital tools for my work. & मुझे डिजिटल उपकरण उपयोग करने की अपनी क्षमता पर विश्वास है। \\ \hline
Q19 & Learning to use a new portal/app takes me: & एक नया पोर्टल/ऐप सीखने में मुझे लगता है: \\ \hline
Q20 & The design/interface of most government portals is user-friendly. & अधिकांश सरकारी पोर्टलों का डिज़ाइन उपयोगकर्ता-अनुकूल है। \\ \hline
Q21 & My superiors encourage the use of digital tools. & मेरे वरिष्ठ अधिकारी डिजिटल उपकरणों के उपयोग को प्रोत्साहित करते हैं। \\ \hline
Q22 & My colleagues regularly use digital tools in their work. & मेरे सहकर्मी नियमित रूप से डिजिटल उपकरणों का उपयोग करते हैं। \\ \hline
Q23 & There is a formal mandate/order requiring digital tool use in my department. & मेरे विभाग में डिजिटल उपकरण के लिए औपचारिक आदेश है। \\ \hline
Q24 & Which general tools are you aware of? & आप किन सामान्य उपकरणों से अवगत हैं? \\ \hline
Q25 & How often do you use digital tools for work? & आप कितनी बार काम के लिए डिजिटल उपकरणों का उपयोग करते हैं? \\ \hline
Q26 & What percentage of your work is done digitally? & आपके कार्य का कितना प्रतिशत डिजिटल रूप से होता है? \\ \hline
Q27 & Have you attended any digital skills training in the last 2 years? & क्या पिछले 2 वर्षों में आपने कोई डिजिटल कौशल प्रशिक्षण लिया है? \\ \hline
Q28 & If yes, how many training sessions? & यदि हाँ, तो कितने प्रशिक्षण सत्र? \\ \hline
Q29 & Rate the quality of training received: & प्राप्त प्रशिक्षण की गुणवत्ता का मूल्यांकन करें: \\ \hline
Q30 & Was the training sufficient for your actual job needs? & क्या प्रशिक्षण आपकी वास्तविक कार्य आवश्यकताओं के लिए पर्याप्त था? \\ \hline
Q31 & What topics need more training? & किन विषयों में अधिक प्रशिक्षण चाहिए? \\ \hline
Q32 & What issues do you face? (tick all) & आपको क्या समस्याएँ आती हैं? (सभी पर टिक करें) \\ \hline
Q33 & How often do digital issues disrupt your work? & डिजिटल समस्याएँ कितनी बार आपके काम में बाधा डालती हैं? \\ \hline
Q34 & I feel comfortable asking for help with digital tools. & मुझे डिजिटल उपकरणों के लिए सहायता माँगने में सहजता है। \\ \hline
Q35 & My organisation provides adequate support for digital tools. & मेरा संगठन डिजिटल उपकरणों के लिए पर्याप्त सहायता प्रदान करता है। \\ \hline
Q36 & My department is committed to digital transformation. & मेरा विभाग डिजिटल परिवर्तन के प्रति प्रतिबद्ध है। \\ \hline
Q37 & Rank these priorities (1=highest, 5=lowest): & इन प्राथमिकताओं को क्रम दें (1=सर्वोच्च, 5=न्यूनतम): \\ \hline
Q38 & What one change would most improve your use of digital tools? & कौन-सा एक बदलाव आपके डिजिटल उपकरण उपयोग में सबसे अधिक सुधार करेगा? \\ \hline
Q39 & Has digital tool adoption improved service delivery to citizens? & क्या डिजिटल उपकरणों से नागरिकों को सेवा वितरण में सुधार हुआ है? \\ \hline
\end{longtable}

\subsection{Role-Specific Questions}
\begin{longtable}{|p{6cm}|p{6cm}|}
\hline
\textbf{English Question} & \textbf{Hindi Translation} \\ \hline
\endfirsthead
\hline
\textbf{English Question} & \textbf{Hindi Translation} \\ \hline
\endhead
\hline
\endfoot
\hline
\endlastfoot
Primary interaction with digital tools: & डिजिटल उपकरणों के साथ आपकी प्राथमिक भूमिका: \\ \hline
Is there one person who does most portal work for others? & क्या कोई एक व्यक्ति दूसरों का पोर्टल कार्य करता है? \\ \hline
Is the training appropriate for your specific job role? & क्या प्रशिक्षण आपकी विशिष्ट भूमिका के लिए उपयुक्त है? \\ \hline
When a portal gives an error, what do you do? & जब पोर्टल त्रुटि देता है, तो आप क्या करते हैं? \\ \hline
Do senior officers use digital tools themselves? & क्या वरिष्ठ अधिकारी स्वयं डिजिटल उपकरण उपयोग करते हैं? \\ \hline
Have digital tools changed the work expected at your level? & क्या डिजिटल उपकरणों ने आपके स्तर पर अपेक्षित कार्य बदला है? \\ \hline
Are there digital tasks beyond your current skill level? & क्या कोई डिजिटल कार्य आपके कौशल स्तर से परे हैं? \\ \hline
Should training differ based on job level? & क्या प्रशिक्षण पद स्तर के अनुसार अलग होना चाहिए? \\ \hline
\end{longtable}

\subsection{Personal Tools Questions}
\begin{longtable}{|p{6cm}|p{6cm}|}
\hline
\textbf{English Question} & \textbf{Hindi Translation} \\ \hline
\endfirsthead
\hline
\textbf{English Question} & \textbf{Hindi Translation} \\ \hline
\endhead
\hline
\endfoot
\hline
\endlastfoot
Do you use any non-government apps or tools to help with your work? & क्या आप अपने काम में कोई गैर-सरकारी ऐप्स या उपकरण उपयोग करते हैं? \\ \hline
Which of these do you use for work-related tasks? (tick all) & इनमें से कौन-से आप कार्य के लिए उपयोग करते हैं? (सभी पर टिक करें) \\ \hline
What do you mainly use these tools for? (tick all) & आप इन उपकरणों का मुख्य रूप से किसलिए उपयोग करते हैं? (सभी पर टिक करें) \\ \hline
How often do you use these personal tools for official work? & आप कितनी बार इन व्यक्तिगत उपकरणों का आधिकारिक काम के लिए उपयोग करते हैं? \\ \hline
Do you feel these tools fill a gap that government systems don't cover? & क्या ये उपकरण सरकारी प्रणालियों की कमी पूरी करते हैं? \\ \hline
Any concerns about using personal tools for official work? & व्यक्तिगत उपकरणों के आधिकारिक उपयोग में कोई चिंता? \\ \hline
\end{longtable}

\end{document}
